\documentclass[a4paper,11pt,fleqn]{article}%fleqn pour aligner les equations à gauche
\usepackage{/home/rweye/Nextcloud/Documents/Overleaf/Styles/Exercice}


\newcommand{\chnum}{Chapitre 5}
\newcommand{\chtitle}{Cortège électronique}
\newcommand{\dtype}{Exercices b}

\begin{document}
\maketitle


\section{Justifier un nombre de doublets non liants}
Des schémas de Lewis de différentes molécules dont présentés ci-dessous.\\

\begin{center}
\chemfig{H-C(-[2]H)(-[-2]H)-C(=[2]O)-H} \hspace{2cm} \chemfig{\charge{90=\| ,-90=\|, 180=\|}{Cl}-C(=[2]\charge{135=\| ,45=\|}{O})-\charge{90=\| ,-90=\|, 0=\|}{Cl}}\\
\end{center}

\textit{Justifier le nombre de doublets non liants sur les atomes d'oxygène et de chlore.}

\section{Prévoir le nombre de doublets non liants}
Des schémas de Lewis incomplets de différentes molécules sont présentées ci-dessous.\\
\begin{tabular}{>{\centering\arraybackslash}m{.3\textwidth} >{\centering\arraybackslash}m{.4\textwidth}  >{\centering\arraybackslash}m{.25\textwidth}}
\chemfig{H-C(-[2]H)(-[-2]H)-C(-[2]H)(-[-2]H)-N(-[2]H)-H} &
\chemfig{H-O-C(-[2]H)(-[-2]H)-C(-[2]H)(-[-2]H)-O-H} &
\chemfig{H-C(=[2]O)-H}  \rule[2.5cm]{0cm}{0cm} \\
Ethanamine & Ethan-1,2-diol & Méthanal\\
\end{tabular}


\textit{Recopier ces schémas de Lewis incomplets, puis les compléter en ajoutant un (ou des) doublet(s) non liant(s). Justifier}

\section{A propos de l'oxyde de baryum}
L'oxyde de baryum BaO est utilisé en tant qu'additif dans la synthèse de verre Crown. Ce type de verre permet d'optimiser les systèmes optiques comme les télescopes.\\
Données :
Configuration électronique d'un atome d'oxygène O : $1s^2 2s^2 2p^4$.\\
L'oxyde de baryum BaO contient des ions monoatomiques issus des atomes des éléments oxygène O et de baryum Ba.\\
\textit{Déterminer à quelle colonne du tableau périodique appartient l'élément baryum Ba.}

\section{Schéma de Lewis de l'acide cyanhydride}
L'acide cyanhydride HCN est une substance toxique que l'on trouve dans certains noyaux (pêche, prune, etc.) ou dans les amandes amères. Deux schémas sont donnés ci-dessous.\\
\begin{center}

\begin{tabular}{|c|c|}
\hline
Proposition 1 & Proposition 2\\
\hline
\chemfig{H-\charge{90=\| ,-90=\|}{C}-\charge{90=\| ,-90=\|, 0=\|}{N}} & \chemfig{H-C~\charge{0=\| }{N}} \rule[.5cm]{0cm}{0cm}\\
\hline
\end{tabular}
\end{center}
\begin{enumerate}
\item Pour les deux propositions, vérifier que tous les atomes respectent la règle de stabilité.
\item L'élément azote se trouve à la 2$^e$ période et 15$^e$ colonne du tableau périodique. Dénombrer les électrons de valence.
\item Déterminer le nombre d'électrons qui appartient "en propre" à l'atome d'azote dans chacune des propositions.
\item En déduire le schéma de Lewis correspondant à l'acide cyanhydride HCN.
\end{enumerate}

\section{L'odeur du gaz de ville}
Le méthane utilisé comme "gaz de ville" est un combustible incolore et inodore. Pour des raisons de sécurité, il est odorisé à l'aide d'une espèce chimique, l'éthylmercaptan, dont la formule brute est $C_2 H_6 S$. Lors d'une fuite de gaz, l'odeur soufrée de l'éthylmercaptan est rapidement perçue.\\

Document A : Schémas de Lewis de molécules contenant l'élément oxygène O\\
\begin{tabular}{>{\centering\arraybackslash}m{.3\textwidth} >{\centering\arraybackslash}m{.4\textwidth}  >{\centering\arraybackslash}m{.25\textwidth}}
\chemfig{H-C(-[2]H)(-[-2]H)-\charge{90=\| ,-90=\|}{O}-C(-[2]H)(-[-2]H)-H} &
\chemfig{H-C(-[2]H)(-[-2]H)-C(-[2]H)(-[-2]H)-\charge{90=\| ,-90=\|}{O}-H} & \chemfig{H-\charge{90=\| ,-90=\|}{O}-H}\\
Méthoxyméthane & Ethanol & Eau\\

\end{tabular}

\vspace{.5cm}
Document B : Schémas de Lewis de molécules contenant des éléments d'une même famille (carbone et silicium)\\
\begin{tabular}{>{\centering\arraybackslash}m{.3\textwidth} >{\centering\arraybackslash}m{.4\textwidth}}
\chemfig{H-C(-[2]H)(-[-2]H)-H} & \chemfig{H-Si(-[2]H)(-[-2]H)-H}\\
Méthane & Silane\\
\end{tabular}

\vspace{.5cm}

\textit{Proposer un schéma de Lewis de la  molécule d'éthylmercaptan.}

\section{Énergie de liaison double}
Voici deux schémas de Lewis incomplets de molécules contenant des liaisons doubles.\\

\vspace{.5cm}

\begin{tabular}{>{\centering\arraybackslash}m{.3\textwidth} >{\centering\arraybackslash}m{.4\textwidth}}
\chemfig{H-C(=[2]O)-H} & \chemfig{H-C(=[2]N-H)-H}\\
Méthanal & Méthanaimine\\
\end{tabular}

\vspace{.5cm}


\begin{enumerate}
	\item Recopier les schémas de Lewis incomplets, puis les compléter en ajoutant un (ou des) doublet(s) non liant(s), sachant que chaque atome vérifie la règle de stabilité.
	\item L'énergie d'atomisation est l'énergie à fournir pou rompre toutes les liaisons d'une molécule et obtenir des atomes :\\ a) Calculer les énergies de liaisons des liaisons C=O et C=N.\\ b) En déduire, parmi les doubles liaisons C=O et C=N, celle qui est la plus stable.
\end{enumerate}

Données :\\
$E_\text{liaison} (C-H) = 413$ U.S.I\\
$E_\text{liaison} (N-H) = 391$ U.S.I\\
$E_\text{atomisation} (\text{méthanal}) = 1567$ U.S.I\\
$E_\text{atomisation} (\text{méthanomine}) = 1564$ U.S.I\\
\end{document}